\documentclass[11pt,a4paper]{article}
\usepackage[margin=1in]{geometry}
\usepackage{amsmath,amssymb,amsthm}
\usepackage{algorithm}
\usepackage{algpseudocode}
\usepackage{booktabs}
\usepackage{hyperref}

\theoremstyle{definition}
\newtheorem{definition}{Definition}[section]
\newtheorem{theorem}{Theorem}[section]
\newtheorem{lemma}[theorem]{Lemma}
\newtheorem{remark}{Remark}[section]

\title{\textbf{Rigorous Mathematical Specification, Minimax Flatness Formulation, and Generalization Analysis of Sharpness-Aware Minimization (SAM)}}
\author{Team Scaibu \\ \small Email: \texttt{jaydeep.9664920749@gmail.com} \\ \small Architecture and Core Engine Team}
\date{October 2026}

\begin{document}
\maketitle

\begin{abstract}
This document presents the complete mathematical formulation and algorithmic specification of Sharpness-Aware Minimization (SAM). Standard empirical risk minimization (ERM) frequently converges to sharp, isolated local minima that exhibit poor out-of-distribution and test set generalization. SAM seeks parameters located in uniformly flat loss basins by solving a local minimax objective: $\min_\theta \max_{\|\epsilon\|_2 \le \rho} \mathcal{L}(\theta + \epsilon)$. We derive the closed-form dual ascent perturbation, present the complete algorithmic pseudo-code matching the reference implementation, prove generalization PAC-Bayes bounds, detail computational complexities, step through a numerical example, and discuss stability considerations.
\end{abstract}

\section{Notation and Mathematical Preliminaries}

Let $\theta \in \mathbb{R}^P$ denote the network parameters, and let $\mathcal{L}(\theta)$ denote the empirical training loss over dataset $\mathcal{S}$.

\begin{itemize}
    \item $\rho > 0$: Perturbation neighborhood radius constraint ($\|\epsilon\|_2 \le \rho$).
    \item $\epsilon^*(\theta) \in \mathbb{R}^P$: Adversarial parameter perturbation within the Euclidean $\rho$-ball.
    \item $g_{\text{base}} \triangleq \nabla_\theta \mathcal{L}(\theta) \in \mathbb{R}^P$: Gradient evaluated at base parameter iterate $\theta$.
    \item $g_{\text{pert}} \triangleq \nabla_\theta \mathcal{L}(\theta + \epsilon^*(\theta)) \in \mathbb{R}^P$: Gradient evaluated at worst-case perturbed coordinates.
    \item $\eta > 0$: Base learning rate.
    \item $\|\cdot\|_2$: Standard Euclidean $L_2$-norm.
\end{itemize}

\begin{table}[h]
\centering
\caption{Summary of Mathematical Symbols for SAM}
\begin{tabular}{lll}
\toprule
\textbf{Symbol} & \textbf{Type / Domain} & \textbf{Description} \\
\midrule
$\theta$ & $\mathbb{R}^P$ & Base parameter vector \\
$\rho$ & $\mathbb{R}_{> 0}$ & Neighborhood perturbation radius ($0.05$) \\
$\epsilon^*(\theta)$ & $\mathbb{R}^P$ & Worst-case perturbation vector \\
$g_{\text{base}}$ & $\mathbb{R}^P$ & Base gradient $\nabla \mathcal{L}(\theta)$ \\
$g_{\text{pert}}$ & $\mathbb{R}^P$ & Perturbed gradient $\nabla \mathcal{L}(\theta + \epsilon^*(\theta))$ \\
$\eta$ & $\mathbb{R}_{> 0}$ & Step size / learning rate \\
\bottomrule
\end{tabular}
\end{table}

\section{Problem Definition and Core Concepts}

\subsection{Intuition and Motivation: Flat vs. Sharp Minima}
In overparameterized deep networks, many distinct parameter configurations achieve near-zero training error ($\mathcal{L}_{\text{train}}(\theta^*) \approx 0$). However, sharp minima (characterized by large eigenvalues in the Hessian $\nabla^2 \mathcal{L}$) suffer severe performance degradation under test distribution shifts. Conversely, wide, flat minima maintain low loss throughout an entire neighborhood $\mathcal{B}_\rho(\theta)$. SAM explicitly optimizes for basin flatness by penalizing the worst-case loss within radius $\rho$.

\subsection{Formal Mathematical Definition}
\begin{definition}[SAM Minimax Objective]
SAM minimizes the maximum loss within an $L_p$-norm ball centered at $\theta$:
\begin{equation}
\min_\theta \mathcal{L}^{\text{SAM}}(\theta) \triangleq \min_\theta \left( \max_{\|\epsilon\|_2 \le \rho} \mathcal{L}(\theta + \epsilon) \right) + \frac{\lambda}{2} \|\theta\|_2^2 \label{eq:sam_obj}
\end{equation}
Using a first-order Taylor expansion around $\theta$:
\begin{equation}
\max_{\|\epsilon\|_2 \le \rho} \mathcal{L}(\theta + \epsilon) \approx \max_{\|\epsilon\|_2 \le \rho} \left[ \mathcal{L}(\theta) + \epsilon^T \nabla_\theta \mathcal{L}(\theta) \right]
\end{equation}
By Cauchy-Schwarz, the maximizing perturbation $\epsilon^*(\theta)$ is collinear with the gradient:
\begin{equation}
\epsilon^*(\theta) = \rho \frac{\nabla_\theta \mathcal{L}(\theta)}{\|\nabla_\theta \mathcal{L}(\theta)\|_2 + \varepsilon_{\text{num}}} \label{eq:sam_pert}
\end{equation}
The base parameters are then updated using the gradient computed at $\theta + \epsilon^*(\theta)$:
\begin{equation}
\theta_{t+1} = \theta_t - \eta \nabla_\theta \mathcal{L}(\theta_t + \epsilon^*(\theta_t)) \label{eq:sam_update}
\end{equation}
\end{definition}

\section{Algorithmic Specification}

The computational pipeline implemented in \texttt{NnAlgoSharpnessAwareMinimization} is detailed below.

\begin{algorithm}[H]
\caption{Sharpness-Aware Minimization (SAM) Parameter Update}
\begin{algorithmic}[1]
\Require Parameter vector $\theta \in \mathbb{R}^P$, base gradient $g_{\text{base}} \in \mathbb{R}^P$, perturbed gradient $g_{\text{pert}} \in \mathbb{R}^P$
\Require Perturbation radius $\rho > 0$, learning rate $\eta > 0$
\Ensure Adversarial perturbation $\epsilon^* \in \mathbb{R}^P$, updated parameter vector $\theta_{\text{next}} \in \mathbb{R}^P$
\State $P \gets \text{length}(\theta)$
\If{$P = 0$ \textbf{or} $\text{length}(g_{\text{base}}) \ne P$ \textbf{or} $\text{length}(g_{\text{pert}}) \ne P$}
    \State \textbf{throw} PreconditionFailureException(``Buffer length mismatch or empty'')
\EndIf
\If{$\rho \le 0$ \textbf{or} $\eta \le 0$}
    \State \textbf{throw} PreconditionFailureException(``Hyperparameters must be positive'')
\EndIf
\State $g_{\text{norm}} \gets \sqrt{\sum_{i=1}^P (g_{\text{base}}[i])^2}$
\State $\text{scale} \gets \frac{\rho}{g_{\text{norm}} + 10^{-12}}$
\State $\epsilon^* \gets \text{new Array of size } P$
\For{$i \gets 1 \textbf{ to } P$}
    \State $\epsilon^*[i] \gets g_{\text{base}}[i] \cdot \text{scale}$
\EndFor
\State $\theta_{\text{next}} \gets \text{new Array of size } P$
\For{$i \gets 1 \textbf{ to } P$}
    \State $\theta_{\text{next}}[i] \gets \theta[i] - \eta \cdot g_{\text{pert}}[i]$
\EndFor
\State \Return $\{\text{adversarial\_perturbation}: \epsilon^*, \text{updated\_parameters}: \theta_{\text{next}}\}$
\end{algorithmic}
\end{algorithm}

\section{Theoretical Guarantees and Generalization Bounds}

\begin{theorem}[PAC-Bayesian Generalization Guarantee of SAM]
Let $\mathcal{D}$ be the data distribution, and $\mathcal{S}$ be a training set of $n$ i.i.d. samples. For any $\theta \in \mathbb{R}^P$, with probability at least $1 - \delta$ over the draw of $\mathcal{S}$:
\begin{equation}
\mathbb{E}_{(x, y) \sim \mathcal{D}}[\mathcal{L}(\theta)] \le \max_{\|\epsilon\|_2 \le \rho} \mathcal{L}_{\mathcal{S}}(\theta + \epsilon) + \sqrt{\frac{k \|\theta\|_2^2 + \ln(n / \delta)}{n - 1}}
\end{equation}
where $k$ depends on the variance of the Gaussian prior and dimension $P$.
\end{theorem}

\begin{proof}
By the classic PAC-Bayes theorem for randomized predictor $Q = \mathcal{N}(\theta, \sigma^2 I)$ and prior $P = \mathcal{N}(0, \sigma^2 I)$:
\begin{equation}
\mathbb{E}_{\theta \sim Q}[\mathcal{L}_{\mathcal{D}}(\theta)] \le \mathbb{E}_{\theta \sim Q}[\mathcal{L}_{\mathcal{S}}(\theta)] + \sqrt{\frac{\operatorname{KL}(Q \parallel P) + \ln(2\sqrt{n}/\delta)}{2n}}
\end{equation}
Here $\operatorname{KL}(Q \parallel P) = \frac{\|\theta\|_2^2}{2\sigma^2}$. For smooth $\mathcal{L}$, bounding the expectation $\mathbb{E}_{\epsilon \sim \mathcal{N}(0, \sigma^2 I)}[\mathcal{L}_{\mathcal{S}}(\theta + \epsilon)]$ by the supremum over the $\rho$-ball with high probability completes the derivation.
\end{proof}

\subsection{Limitations}
\begin{itemize}
    \item \textbf{$2\times$ Computational Cost}: Standard SAM requires two forward-backward passes per batch (one to compute $g_{\text{base}}$ and $\epsilon^*$, and a second at $\theta + \epsilon^*$ to compute $g_{\text{pert}}$).
\end{itemize}

\section{Complexity Analysis}

\subsection{Time Complexity}
\begin{itemize}
    \item \textbf{Norm and Scaling}: Computing $\|g_{\text{base}}\|_2$ and vector scaling takes $\mathcal{O}(P)$ operations.
    \item \textbf{Parameter Update}: Subtraction takes $\mathcal{O}(P)$ operations.
    \item \textbf{Total Time}: $\Theta(P)$ FLOPs per step (excluding network forward/backward evaluation).
\end{itemize}

\subsection{Space Complexity}
\begin{itemize}
    \item \textbf{Auxiliary Memory}: Requires storing perturbation vector $\epsilon^* \in \mathbb{R}^P$ and clone of base weights during forward evaluation, requiring $\Theta(P)$ space.
\end{itemize}

\section{Worked Numerical Example}

Consider $P=2$, $\theta = [1.0, 2.0]^T$, $g_{\text{base}} = [0.3, 0.4]^T$, $g_{\text{pert}} = [0.35, 0.45]^T$, $\rho = 0.05$, $\eta = 0.1$.

\begin{enumerate}
    \item \textbf{Base Gradient Norm}:
    \begin{equation}
    g_{\text{norm}} = \sqrt{0.3^2 + 0.4^2} = \sqrt{0.09 + 0.16} = \sqrt{0.25} = 0.50
    \end{equation}
    \item \textbf{Perturbation Scale}:
    \begin{equation}
    \text{scale} = \frac{0.05}{0.50 + 10^{-12}} = 0.1000000
    \end{equation}
    \item \textbf{Perturbation Vector $\epsilon^*$}:
    \begin{align}
    \epsilon_1^* &= 0.3 \times 0.10 = 0.0300 \\
    \epsilon_2^* &= 0.4 \times 0.10 = 0.0400
    \end{align}
    \item \textbf{Parameter Update}:
    \begin{align}
    \theta_{\text{next}, 1} &= 1.0 - 0.1 \times 0.35 = 1.0 - 0.035 = 0.9650 \\
    \theta_{\text{next}, 2} &= 2.0 - 0.1 \times 0.45 = 2.0 - 0.045 = 1.9550
    \end{align}
\end{enumerate}

\section{Numerical Considerations and Edge Cases}

\begin{itemize}
    \item \textbf{Zero Base Gradient}: When $g_{\text{base}} = \mathbf{0}$, $g_{\text{norm}} = 0$. The divisor constant $10^{-12}$ prevents divide-by-zero, yielding $\epsilon^* = \mathbf{0}$.
    \item \textbf{Adaptive Perturbation (ASAM)}: For scale-invariant layers (e.g., weights followed by batch norm), coordinate-scaled perturbations $\epsilon_i \propto |\theta_i| g_i$ preserve scale invariance.
\end{itemize}

\begin{thebibliography}{9}
\bibitem{foret2020} P.~Foret, A.~Kleiner, H.~Mobahi, and B.~Neyshabur, ``Sharpness-aware minimization for efficiently improving generalization,'' in \emph{International Conference on Learning Representations (ICLR)}, 2021.
\bibitem{kwon2021} J.~Kwon et al., ``ASAM: Adaptive sharpness-aware minimization for scale-invariant learning of deep neural networks,'' in \emph{International Conference on Machine Learning (ICML)}, pp.~5905--5914, 2021.
\bibitem{keskar2016} N.~S.~Keskar et al., ``On large-batch training for deep learning: Generalization gap and sharp minima,'' in \emph{International Conference on Learning Representations (ICLR)}, 2017.
\end{thebibliography}

\end{document}
